\bookchapter{Bab 19: Pengesanan Kerosakan \& Penyelamatan Diri}

\chapterintro{Periksa ruang penyimpanan}

\booksection{: Pergi ke Tetapan Penyimpanan lihat ruang yang tinggal. Jika kurang daripada 10\%, sistem akan menjadi sangat lengah, kerana ia berjuang untuk membebaskan ruang. Pengalaman saya: Kekalkan sekurang-kurangnya 15\%-20\% ruang yang tinggal, telefon pintar baru boleh berjalan dengan agak lancar.}

Periksa kedudukan penggunaan bateri

\booksubsection{: Pergi ke tetapan bateri lihat kedudukan penggunaan bateri dalam 24 jam lepas. Jika satu App menduduki bahagian besar tetapi anda hampir tidak menggunakannya, itu bermakna ia berjalan di latar belakang secara senyap. Nyahpasang terus atau hadkan aktiviti latar belakangnya.}

\begin{enumerate}
  \item \textbf{Fikirkan apa yang baru dilakukan}: Adakah ia menjadi lengah selepas menaik taraf sistem? (Biasa menjadi lengah beberapa hari pertama selepas kemas kini versi besar, sistem membina semula indeks di latar belakang). Atau adakah ia menjadi lengah selepas memasang App baru? (Cuba nyahpasang).
  \item \textbf{Lihat persekitaran penggunaan}: Telefon pintar terdedah kepada matahari, memakai sarung tebal sambil mengecas dan bermain permainan, kepanasan adalah pasti. Suhu tinggi juga menyebabkan CPU secara automatik mengurangkan frekuensi, anda akan rasa telefon pintar tiba-tiba menjadi lengah. Biarkan sejuk dan cuba lagi.
  \item \textbf{Kaedah rawatan (dari mudah ke kompleks)}: Mulakan semula telefon pintar (menyelesaikan 60\% masalah) Kosongkan ruang penyimpanan (padam App dan video yang tidak digunakan) Jika satu App menyebabkan masalah, kosongkan cache atau nyahpasang dan pasang semula Langkah terakhir: Sandar data dan pulihkan tetapan kilang.
  \item \textbf{2. App Keluar Tiba-tiba, Tidak Responsif — App Tidak Mendengar Arahan}Satu App keluar serta-merta apabila dibuka, atau tiba-tiba kaku dan tidak bergerak semasa digunakan. Proses pengendalian seperti doktor membuat diagnosis, langkah demi langkah:
\end{enumerate}

\textbf{Tutup dan buka semula}: Tatal App dari latar belakang, kemudian buka semula. Kebanyakan pepijat sementara akan selesai dengan langkah ini.

\booksection{Kemas kini App}

: Pergi ke kedai aplikasi lihat ada kemas kini atau tidak. Kadang-kadang kerana versi App terlalu lama, tidak serasi dengan sistem.

\begin{enumerate}
  \item \textbf{Periksa kebenaran}: Sesetengah App memerlukan kebenaran tertentu untuk berfungsi dengan normal, tetapi anda tidak memberikannya. Contohnya App kamera tidak diberikan kebenaran kamera, dibuka terus keluar. Pergi ke tetapan cari App tersebut, berikan semua kebenaran yang sepatutnya.
  \item \textbf{Kosongkan cache}: Pergi ke Tetapan Pengurusan Aplikasi cari App tersebut Penyimpanan Kosongkan Cache. Perhatian: "Kosongkan Cache" bukan "Kosongkan Data"! Mengosongkan data akan memadam status log masuk dan tetapan dalam App anda, bersamaan dengan memulihkan tetapan kilang untuk App tersebut.
  \item \textbf{Nyahpasang dan pasang semula}: Langkah terakhir. Sebelum nyahpasang, pastikan data App ini tidak akan hilang — untuk WeChat, sandar rekod sembang terlebih dahulu; untuk permainan, pastikan akaun telah diikat dengan nombor telefon.
  \item \textbf{3. Tidak Boleh Dicas atau Cas Perlahan — Mungkin Bukan Salah Bateri}Tidak boleh dicas adalah masalah yang sangat mudah menyebabkan kebimbangan, tetapi kebanyakan kali punca adalah pada aksesori, bukan telefon pintar rosak.
  \item \textbf{Periksa port pengecasan terlebih dahulu}: Gunakan pic gigi atau alat kecil bukan konduktor, perlahan-lahan membersihkan habuk dan serat kapas dalam port pengecasan. Saya sendiri pernah mengalaminya — seutas benang dari poket tersangkut di port pengecasan, apabila dimasukkan disangka sentuhan tidak baik, selepas dikeluarkan ia kembali normal.
\end{enumerate}

\booksection{Cuba tukar kabel dan pengecas}

: Kadar kerosakan kabel jauh lebih tinggi daripada pengecas, terutamanya kabel yang sering digulung. Tukar satu kabel dan pengecas yang pasti tiada masalah untuk ujian bersilang, cepat boleh tentukan aksesori mana yang rosak.

\begin{itemize}
  \item \textbf{Masalah pengecasan wayarles}: Jika anda menggunakan pengecasan wayarles, periksa sama ada sarung telefon pintar terlalu tebal, atau ada kepingan logam di belakang sarung (seperti dudukan magnet). Juga, letakkan telefon pintar tepat di tengah pad pengecasan — jika senget mungkin tidak boleh dicas.
  \item \textbf{Cas perlahan}: Sahkan kuasa pengecas mencukupi atau tidak (iPhone mahu cas pantas memerlukan sekurang-kurangnya pengecas PD 20W), sama ada kabel menyokong cas pantas (bukan semua kabel boleh melepasi arus melebihi 3A), sama ada telefon pintar terlalu panas (terlalu panas akan mengurangkan kuasa secara automatik untuk melindungi bateri).
  \item \textbf{4. Rangkaian Tidak Boleh Disambung, Wi-Fi Sentiasa Terputus}Di luar menggunakan data baik-baik saja, pulang ke rumah tidak boleh menyambung Wi-Fi; atau Wi-Fi disambung tetapi tiada internet. Sahkan mengikut susunan ini:
  \item \textbf{Hidap/matikan mod penerbangan}: Tarik ke bawah pusat kawalan, hidupkan mod penerbangan tunggu 5 saat kemudian matikan. Ini bersamaan dengan membiarkan telefon pintar mencari rangkaian semula, sangat berkesan.
\end{itemize}

\booksection{“Lupakan” Wi-Fi dan sambung semula}

: Pergi ke tetapan Wi-Fi, ketik tanda seru di sebelah nama rangkaian tersebut (Android ialah gear), pilih “Abai Rangkaian Ini” atau “Lupakan”, kemudian masukkan kata laluan semula untuk menyambung.

\begin{enumerate}
  \item \textbf{Mulakan semula penghala}: Ramai orang di rumah tidak pernah mematikan penghala selama berbulan-bulan, lama-kelamaan ia menjadi tidak stabil. Cabut kuasa tunggu 30 saat kemudian pasang semula, kebanyakan masalah Wi-Fi akan selesai.
  \item \textbf{Periksa sama ada ia masalah telefon pintar}: Adakah peranti lain menyambung Wi-Fi yang sama dengan normal? Jika semuanya normal hanya telefon pintar anda bermasalah, telefon pintar mungkin perlu dimulakan semula atau tetapan semula rangkaian (Tetapan Umum Pindah atau Tetap Semula Tetap Semula Tetapan Rangkaian, tidak akan memadam data anda, tetapi akan mengosongkan semua kata laluan Wi-Fi).
  \item \textbf{5. Masalah Skrin: Sentuhan Tidak Sensitif, Skrin Berkelip, Meninggalkan Jejak}Sentuhan tidak sensitif
  \item \textbf{: Tanggalkan pelindung kaca terlebih dahulu lihat sama ada salah pelindung — pelindung kaca murahan memang boleh menjejaskan sentuhan selepas lama digunakan. Terdapat air pada skrin atau minyak pada tangan juga boleh menyebabkan sentuhan melayang. Jika bukan kedua-duanya, mulakan semula dan cuba lagi.}Jejak/bayangan pada skrin (terbakar)
\end{enumerate}

\booksection{: Memaparkan gambar yang sama untuk masa yang lama (contohnya sentiasa menyala menatal TikTok), akan meninggalkan “bayangan” pada skrin OLED. Ini bukan skrin rosak, hanya sementara. Matikan skrin rehat setengah jam biasanya boleh pulih.}

\begin{itemize}
  \item \textbf{Skrin berkelip}: Jika berkelip pada kecerahan rendah, mungkin masalah modulasi PWM, biasa pada skrin OLED. Pergi ke tetapan cari pilihan "Pengmodulan DC" atau "Anti-kelip", hidupkan akan lebih baik. Tidak semua telefon pintar mempunyai fungsi ini.
  \item \textbf{Skrin ada satu kawasan yang hitam/terang sepenuhnya}: Ini agak serius, kebanyakannya kerosakan fizikal skrin, perlu menukar skrin. Skrin yang diganti bukan secara rasmi, warna dan kecerahan biasanya tidak sebaik ori, dan juga menjejaskan kalis air.
  \item \textbf{6. Tidak Boleh Dihidupkan, Tersekat di Antara Muka Logo — Jangan Panik Dulu}Situasi ini paling mudah membuatkan degupan jantung meningkat, tetapi bertenang, cuba ini dulu:
  \item \textbf{Cas setengah jam dahulu}: Kebanyakan kali hanya bateri benar-benar habis. Pasang pengecas ori, tunggu sekurang-kurangnya 30 minit, kemudian cuba hidupkan. Sesetengah telefon pintar perlu dicas ke tahap tertentu selepas bateri habis baru akan bertindak balas.
\end{itemize}

\booksection{Mula semula secara paksa}

: iPhone skrin penuh tekan sekali volume naik, sekali volume turun, tahan butang kuasa sisi tanpa lepaskan, sehingga Logo Apple muncul pada skrin (mungkin perlu ditahan 10-15 saat). Kaedah universal Android: Tahan butang kuasa lebih 10 saat, sehingga bergetar atau skrin menyala.

\begin{itemize}
  \item \textbf{Tersekat di antara muka permulaan berulang kali mula semula}: Situasi ini dipanggil "Apple Putih" (iPhone) atau "Tersangkut di Logo" (Android). Jika mula semula secara paksa tidak berjaya, mungkin perlu memasuki mod pemulihan untuk membaiki dengan komputer. iPhone sambung ke komputer buka iTunes/Finder, masuk mod pemulihan kemudian pilih "Kemas Kini" (bukan "Pulihkan"), akan memasang semula sistem tetapi mengekalkan data. Android masuk mod Fastboot atau Recovery, kombinasi kekunci setiap jenama berbeza, disyorkan cari model telefon pintar anda + "masuk mod pemulihan".
  \item \textbf{Sambung semula Wi-Fi}: Maka hanya boleh meminta bantuan pembaikan profesional. Jika data telah disandar, tidak perlu terlalu bimbang.
  \item \textbf{7. Masalah Bunyi: Pembesar Suara Tidak Berbunyi, Bunyi Panggilan Kecil}Pembesar suara tiada bunyi
  \item \textbf{: Pastikan bukan mod senyap (tuil senyap di sisi iPhone), mod jangan ganggu, atau volum ditekan ke paling rendah. Kedengaran bodoh, tetapi situasi ini merangkumi 40\%. Kemudian periksa sama ada Bluetooth disambung ke fon kepala atau pembesar suara tertentu.}Pihak lawan tidak mendengar anda bercakap semasa panggilan
\end{itemize}

\booksection{: Periksa sama ada lubang mikrofon tersumbat — lubang kecil di bahagian bawah telefon pintar sangat mudah tersumbat oleh habuk. Bersihkan perlahan-lahan dengan berus lembut. Jangan gunakan jarum atau pic gigi untuk mencucuk keras, akan merosakkan jejaring anti-debu.}

\begin{itemize}
  \item \textbf{Bunyi serak/pecah}: Jika disebabkan oleh pembesar suara kemasukan air — jangan gunakan pengering rambut angin panas! Gunakan tisu kering menyerap kelembapan permukaan, kemudian biarkan kering secara semula jadi. Dalam talian kata App saliran gelombang bunyi ada kesan tertentu (memainkan bunyi frekuensi rendah untuk menggetarkan air keluar), boleh cuba.
  \item \textbf{8. Cari Pembaikan Rasmi, Elakkan Perangkap Tersebut}Jika semua kaedah di atas telah dicuba tetapi masih tidak berjaya, maka perlu mencari pembaikan profesional. Tetapi pasaran pembaikan sangat dalam, berikut adalah pengalaman dan perangkap yang pernah saya langgar sendiri:
  \item \textbf{Cari rasmi sebenar, jangan cari yang "kelihatan seperti rasmi"}: Pergi ke laman sesawang rasmi jenama (Apple.com.cn, Huawei.com.cn dll.) cari senarai "Pembekal Perkhidmatan Sah". Kedai bukan sahih di tepi jalan yang menggantung logo Apple, mungkin menggunakan aksesori bukan ori, selepas pembaikan kalis air anda hilang.
\end{itemize}

\booksection{Perkara yang mesti dilakukan sebelum menghantar pembaikan}

: Sandar semua data (jangan harap data tidak hilang semasa pembaikan), log keluar semua akaun penting (Apple ID, WeChat, Alipay), matikan “Cari Telefon Pintar Saya”. Sesetengah kedai pembaikan akan meminta anda menyediakan kata laluan kunci skrin, ini ada risiko privasi — lebih baik anda sendiri log keluar akaun terlebih dahulu sebelum menghantar.

\begin{itemize}
  \item \textbf{Perangkap menukar skrin}: Skrin bukan rasmi dibahagi kepada "skrin tekan ori" (kaca ori diganti), "skrin tekan belakang" (kaca bukan ori), "skrin dipasang" (semua bukan ori). Perbezaan harga beberapa kali ganda, tetapi kesan paparan dan ketahanan juga sangat berbeza. Jika telefon pintar masih dalam tempoh waranti, ikut rasmi dengan jujur.
  \item \textbf{Perangkap menukar bateri}: Bateri bukan ori mungkin memalsukan kapasiti, dan selepas digunakan setengah tahun penurunannya lebih besar daripada bateri ori yang digunakan selama dua tahun. Harga rasmi menukar bateri Apple tidak jauh lebih mahal daripada bukan rasmi, disyorkan ikut rasmi.
  \item \textbf{Perangkap pemulihan data}: Selepas telefon pintar kemasukan air/jatuh rosak, pergi membaiki jika pihak lain berkata "pemulihan data perlu tambah bayaran", berhati-hati — sesetengah peniaga tidak jujur akan sengaja membesarkan tahap kerosakan untuk mengenakan bayaran lebih. Pemulihan data yang sah akan memetakan harga terlebih dahulu sebelum beroperasi, dan jika pemulihan tidak berjaya tiada caj.
  \item \textbf{Langkah Pertama Universal: Mulakan Semula Telefon Pintar}Tidak kira masalah apa yang anda hadapi — lengah, keluar tiba-tiba, tidak boleh menyambung rangkaian, bunyi luar biasa — mulakan semula telefon pintar terlebih dahulu. Ayat ini kedengaran seperti mengelak, tetapi berdasarkan pengalaman saya sendiri, sekurang-kurangnya 60\% masalah kecil akan selesai selepas dimulakan semula. Telefon pintar seperti komputer, selepas berjalan lama, ingatan akan terkumpul pelbagai benda kotor, mulakan semula ialah mengosongkan dan bermula dari awal. Jadi, kali depan telefon pintar bermasalah, jangan tergesa-gesa cari "apa yang perlu dilakukan", mulakan semula dahulu kemudian lihat.
  \item \textbf{Hafalkan sekarang: Operasi penyelamat ini perlu dihafal di luar kepala}Mula semula secara paksa
\end{itemize}

\begin{tcbwarn}{: iPhone skrin penuh tekan volume naik volume turun tahan butang kuasa sehingga Logo muncul. Android tahan butang kuasa 10-15 saat sehingga bergetar.}
Bersihkan port pengecasan

\end{tcbwarn}

\begin{tcbpractice}{: Cari alat kecil (pic gigi atau pen jarum telefon pintar), perlahan-lahan bersihkan habuk dalam port pengecasan.}
\begin{enumerate}
  \item \textbf{Sambung semula Wi-Fi}: Pergi ke tetapan Wi-Fi "lupakan" Wi-Fi semasa anda, masukkan semula kata laluan untuk menyambung.
  \item \textbf{Hidup/matikan mod penerbangan}: Tarik ke bawah pusat kawalan, hidupkan mod penerbangan, kira 5 saat dalam hati, matikan, lihat sama ada rangkaian dipulihkan.
  \item \textbf{Tangkap layar langkah-langkah operasi ini dan simpan, atau tulis dalam nota. Apabila menghadapi masalah tidak perlu panik mencari, terus buka nota sahaja.}Bab 20: Mod Pembangun \& Kebenaran Lanjutan
  \item \textbf{Bab ini disediakan untuk pengguna yang ingin menerokai "fungsi tersembunyi" telefon pintar dengan lebih mendalam. Mod pembangun, kebenaran Root, kunci BL... konsep-konsep yang kedengaran sangat geek ini sebenarnya tidak jauh dari pengguna biasa.}1. Mod Pembangun: Buka Kunci Tetapan Tersembunyi
  \item Mod pembangun ialah menu tersembunyi yang terbina dalam sistem telefon pintar, terdapat banyak fungsi yang tidak diperlukan oleh pengguna biasa tetapi diperlukan oleh pembangun.
\end{enumerate}

\end{tcbpractice}
